% Generated draft: requires mathematical and source review.
\documentclass[11pt]{article}
\usepackage[margin=25mm]{geometry}
\usepackage{fontspec}
\setmainfont{Latin Modern Roman}
\usepackage{amsmath,amssymb,mathtools,unicode-math}
\setmathfont{Latin Modern Math}
\usepackage{enumitem,needspace,xurl,hyperref}
\setlength{\parindent}{0pt}
\setlength{\parskip}{5pt}
\setlength{\emergencystretch}{3em}
\newcommand{\N}{\mathbb{N}}
\newcommand{\Z}{\mathbb{Z}}
\newcommand{\Q}{\mathbb{Q}}
\newcommand{\R}{\mathbb{R}}
\newcommand{\C}{\mathbb{C}}
\newcommand{\F}{\mathbb{F}}
\newcommand{\A}{\mathbb{A}}
\newcommand{\PP}{\mathbb{P}}
\newcommand{\E}{\mathbb{E}}
\newcommand{\eps}{\varepsilon}
\newcommand{\dd}{\,\mathrm{d}}
\begin{document}
\section*{UnsolvedMath ID: 5000004}
\textbf{Source problem code:} catalogue entry
\subsection{Definitions and mathematical statement}
Fix \(n\ge5\), \(\lambda=2\cos(\pi/n)\), and \(H=\{(x,y)\in\R^2:y>0\}\). A short polygonal billiard trajectory joins vertices and encounters no vertex between them; successive reflections across encountered sides unfold it into a straight segment. Reachable points are endpoint vertices of such unfolded segments, in the source's fixed coordinates. Let \(R_k\) denote those of trajectory type \(A_k\), with indices modulo \(n-2\), and let \(R\) be the full reachable set. The source defines the endpoint angles \(\alpha,\beta\) in its Definition 2.1; the supplied record gives only the type condition \(A_0:\beta-\alpha=2\pi/n\).

Suppose \(u,v\in R_0\) form a unitary pair:
\[
\det(u,v)=u_1v_2-u_2v_1=\sin((n-2)\pi/n).
\]
For \(m\in\Z\) set
\[
u_m=u+m(\lambda+1)v,\qquad
w_m=u+\bigl(\lambda+m(\lambda+1)\bigr)v.
\]
The conjecture asserts
\[
\forall m\in\Z,\quad u_m\in H\Longrightarrow u_m\in R_0,
\qquad w_m\in H\Longrightarrow w_m\in R_1,
\]
and exhausts the reachable points on the line in the relevant upper-half-plane domain:
\[
R\cap H\cap(u+\R v)
=\{u_m:m\in\Z,\ u_m\in H\}\cup
\{w_m:m\in\Z,\ w_m\in H\}.
\]
It further asks for the analogous arithmetic description on \(v+\R u\), where the only reachable types are \(A_0\) and \(A_{n-3}\), the latter being \(A_{-1}\) modulo \(n-2\). In particular it requires
\[
R\cap H\cap(v+\R u)\subseteq R_0\cup R_{n-3}.
\]
The exact analogous coefficients on that second line are not supplied, so a sign or orientation convention is not invented by merely interchanging \(u,v\), which reverses the determinant.

The coordinate origin, polygon scale, angle definitions, and remaining type definitions are indispensable source data missing from the extraction. Also, the wording no other reachable points could mean a full-line statement if reachability is not already restricted to \(H\); the displayed equality states the upper-half-plane version supported by the qualifications in (a)--(b). The supplied research summary reports the arithmetic parametrization and exhaustion claim unresolved, with no later resolution located through 2026.
\subsection{Short English statement}
For a unitary pair of unfolded regular-polygon billiard endpoints, prove the two stated arithmetic families give exactly the reachable points on the first line, with the corresponding two-type description on the second. Source angle, type, and orientation conventions remain necessary.
\subsection{Sources}
\begin{itemize}
\item Local UnsolvedMath registry, permanent ID \(5000004\), exact source code \texttt{catalogue entry}, ``Reachable points on lines through a unitary pair.''
\item Dmitry Fuchs, \emph{Billiard Trajectories in Regular Polygons and Foliations by Closed Geodesics on Surfaces} (2020), Conjecture 2.5 and Definition 2.1; \url{https://amj.math.stonybrook.edu/pdf-Springer-final/020-0170.pdf}. Recorded access 2026-07-29; pdf-text extraction.
\item Supplied OPEN-TRIAGE research summary places this in the conjectural family 2.3--2.7 of that paper, and reports no later resolution located through 2026. No additional reference is supplied.
\end{itemize}
\end{document}
